\section{Model Description}
\label{sec:model}

\subsection{Field Definitions}

The DBS micro-continuum formulation \citep{brinkman1949calculation,soulaine2016micro,
soulaine2024microcontinuum} represents the pore medium as a continuum in which
a cell-centred porosity field $\epsilon$ varies smoothly from
zero (solid) to unity (open pore). The steady incompressible flow is governed
by the Darcy--Brinkman equations:
\begin{align}
\nabla \cdot \bm{u}_D &= 0,
\label{eq:continuity} \\
-\nabla p + \nabla \cdot \bigl(\mu_{\mathrm{eff}} \nabla \bm{u}\bigr)
- \frac{\mu}{K(\epsilon)}\,\bm{u}_D &= \bm{0},
\label{eq:brinkman}
\end{align}
where $\bm{u}_D = \epsilon\bm{u}$ is the Darcy velocity,
$\mu$ is the fluid viscosity, $\mu_{\mathrm{eff}}$ is an effective
viscosity that smoothly recovers the Stokes equation in open-pore cells
\citep{soulaine2016micro}, and $K(\epsilon)$ is a permeability
function. The drag coefficient $\mu/K(\epsilon)$ therefore acts
as a large resistance term in solid-like cells and vanishes in open-pore cells,
so Eq.~\eqref{eq:brinkman} reduces to the Stokes equation.

The reactive-transport model uses two cell-centred material fields. The
porosity $\epsilon$ is the mobile-water volume per bulk cell volume and enters
both DBS transport and the fixed PHREEQC cell mapping. In the analytical
derivations below it is denoted $\theta$ to retain the conventional notation
for mobile-water fraction; thus $\theta=\epsilon$ in every reported case. The
second field is the reactive internal surface area per bulk volume $a_{\mathrm{surf}}$
(m$^2$\,m$^{-3}$), which represents surface available for adsorption and
surface complexation. All aqueous species, together with $\epsilon$,
$a_{\mathrm{surf}}$, and the adsorbed surface fields, are cell-centred; the
transport flux is evaluated on cell faces.

\subsection{Reactive Surface Area Field}
\label{sec:aSurf}

The DBS micro-continuum framework naturally produces three distinct local
material states within a computational cell: a resolved macro-pore (high
$\theta$, negligible grain fraction), a microporous grain (intermediate
$\theta$, significant unresolved grain fraction), and a near-solid or
fully solid region (very low $\theta$). All three must contribute reactive
surface area to the chemistry, but they do so through different geometric
mechanisms. Both contributions are captured in the single cell-centred field
$a_{\mathrm{surf}}$ (units: m$^2\,\mathrm{m}^{-3}_{\mathrm{bulk}}$), which is
the sum of two terms:
\begin{equation}
a_{\mathrm{surf}} = a_{\mathrm{surf}}^{\mathrm{interface}} +
a_{\mathrm{surf}}^{\mathrm{intrinsic}}.
\label{eq:asurf}
\end{equation}

The first term, $a_{\mathrm{surf}}^{\mathrm{interface}}$, is computed
geometrically from the segmented voxel image. The user designates one image
label as the \emph{reference grain phase}; this is simply the segmented label
whose visible image-scale boundary is counted as reactive interface. It may
correspond to a microporous grain, a dense solid, or any other label assigned
to the reactive surface. In general,
$a_{\mathrm{surf}}^{\mathrm{interface}}$ is the local interface area per bulk
cell volume obtained by reconstructing the boundary of that reference phase
inside each finite-volume cell.

The interface contribution is obtained by reconstructing a sub-voxel boundary from the
same segmented image before accumulating area back into the computational
fluid dynamics (CFD) cells. The segmented owner phase is converted to a smoothed
occupancy field, a sub-voxel interface is extracted, and the resulting line or
surface measure is assigned to the owning finite-volume cells as a bulk
$a_{\mathrm{surf}}$ field.
Section~\ref{sec:valSummary} validates this geometric construction separately
from the intrinsic micropore contribution. The resulting total field is
supplied explicitly to the reactive-transport solver.

The interface term captures only grain--fluid boundaries that are visible at
the resolved image scale. It does not detect
reactive area that lies entirely \emph{within} a voxel, such as sub-voxel
surface roughness of a grain or the internal pore network of a microporous
phase. Those surfaces are below the image resolution by construction.

The second term, $a_{\mathrm{surf}}^{\mathrm{intrinsic}}$, accounts for this
sub-voxel contribution. A volumetric surface density $s_k$
(units: m$^2$\,m$^{-3}$) is prescribed for each segmented image label $k$
(for example pore, micropore, or dense solid); it may
represent any sub-voxel surface that the image cannot resolve --- for example
the Brunauer--Emmett--Teller (BET) surface area of a microporous grain measured
by N$_2$ adsorption, or the surface roughness of a dense solid grain measured
on powdered material.
The intrinsic contribution is assembled by weighting each phase density by
its local volume fraction $f_k$ within the coarse cell:
\begin{equation}
a_{\mathrm{surf}}^{\mathrm{intrinsic}} = \sum_k f_k\,s_k.
\label{eq:asurf_intrinsic}
\end{equation}
$s_k$ may be non-zero for any label, including the reference grain phase
itself. The same segmented label may contribute through both terms:
$a_{\mathrm{surf}}^{\mathrm{interface}}$ gives the image-scale outer surface
of the grain, and $a_{\mathrm{surf}}^{\mathrm{intrinsic}}$ adds the
sub-voxel surface area that the image resolution cannot capture. If
$s_k = 0$ for all labels, the total reduces to the purely geometric
image-scale surface. Both contributions are pre-computed from the segmented image
before the flow simulation begins and are held constant throughout the
reactive-transport calculation.

Both contributions enter the geochemical module in the same way. The
PhreeqcRM surface convention used here is based on representative bulk volume.
The volumetric area field is therefore converted only from cubic metres to
litres:
\begin{equation}
a_{\mathrm{surf}}^{\mathrm{bulk,L}} = \frac{a_{\mathrm{surf}}}{10^{3}},
\label{eq:asurf_bulk_litre}
\end{equation}
where the factor $10^3$ converts cubic metres to litres. Surface areas
estimated from nano-computed tomography
(nano-CT), micro-computed tomography (micro-CT), focused ion beam scanning
electron microscopy (FIB-SEM), or other image data are first expressed on a
bulk-volume basis (m$^2$\,m$^{-3}_{\mathrm{bulk}}$), consistent with
Eqs.~\eqref{eq:asurf}--\eqref{eq:asurf_intrinsic}. Image-derived surface areas
reported per microporous-phase volume or per grain mass are converted to this
bulk-volume basis before Eq.~\eqref{eq:asurf_bulk_litre} is applied.

The same fixed porosity field is used for DBS transport and for the PHREEQC
solution volume. Cells in
which $\theta$ falls below the chemistry mask threshold
$\theta_{\mathrm{mask}}$ are treated as chemically inactive and their surface
terms are suppressed.

\subsection{Conservative Aqueous Transport}

For each aqueous species with concentration $C_i$, the conservative DBS
transport equation is written as
\begin{equation}
\frac{\partial (\theta C_i)}{\partial t}
+ \nabla \cdot (\bm{u}_D C_i)
- \nabla \cdot (\theta D_i \nabla C_i) = 0,
\end{equation}
where $\bm{u}_D$ is the Darcy velocity satisfying
Eqs.~\eqref{eq:continuity}--\eqref{eq:brinkman}, and $D_i$ denotes the
diffusivity used in the transport equation. In fully resolved macro-pore cells
($\theta \approx 1$), $D_i$ is the free molecular diffusivity. In cells that
represent microporous material ($0 < \theta < 1$), $D_i$ is interpreted as an
effective diffusivity, commonly $D_i = D_i^{\mathrm{mol}}/\tau$, where
$D_i^{\mathrm{mol}}$ is the molecular diffusivity and $\tau \geq 1$ is the
tortuosity factor of the micropore network. In either case, the $\theta$ factor
multiplying $D_i$ in the diffusion term already accounts for the reduced
water volume available for transport. There is no chemistry source term in
this equation: reactions are handled by the subsequent operator-split
PHREEQC step and do not appear as an explicit right-hand side here.

In the implementation, $\theta$ in this equation is the OpenFOAM field
\texttt{eps}; this field is fixed during all PHREEQC calculations reported here.

For a vector of aqueous species $\bm{C}$, the same transport balance can be
written in
compact form as
\begin{equation}
\frac{\partial (\theta \bm{C})}{\partial t}
+ \nabla \cdot (\bm{u}_D \bm{C})
- \nabla \cdot (\theta \bm{D} \nabla \bm{C}) = \bm{0},
\end{equation}
which makes the conservative structure explicit before any chemistry update is
applied.

\subsection{Cell-Based Surface Inventory}

For the cell-based unresolved surface carried by the DBS formulation, the
conservative local inventory is expressed
in the current solver field variables as
\begin{equation}
M = \theta C + a_{\mathrm{surf}}\Gamma,
\label{eq:inventory}
\end{equation}
where $\Gamma$ is the cell-centred surface-species concentration written by
the implementation (mol\,m$^{-2}$ of reactive surface). The first term is the
aqueous amount per bulk volume and the second is the adsorbed amount per bulk
volume. Porosity enters the aqueous inventory exactly once; the bulk-based
surface-area field requires no additional porosity scaling.

This field-level expression is the appropriate quantity for analytical
comparisons and for benchmark inventory checks. For multiple surface species,
the local bulk inventory is obtained by summing the relevant aqueous and
adsorbed contributions. Single-species notation is used below only for clarity;
all reported totals are evaluated in the actual implemented field variables.

\subsection{Local Chemistry Coupling}

The chemical update is local to each cell. Over one transport--chemistry
cycle, the transport step first advances the aqueous species in conservative
finite-volume form. The updated cell state, together with the mobile-water
fraction and cell-based surface state, is then mapped into PHREEQC
\citep{parkhurst2013phreeqc} via the PhreeqcRM interface
\citep{parkhurst2015phreeqcrm}. PHREEQC solves the local chemistry under the
assumption of \emph{thermodynamic equilibrium}: within each timestep, surface
complexation reactions are treated as instantaneous relative to transport.
This is appropriate when the equilibration timescale of the surface reactions
is short compared with the transport timescale. The relevant transport regime
is characterised by the P\'{e}clet number, $Pe = U L / D$, where $U$ is a
representative velocity, $L$ is a representative transport length, and $D$ is
the aqueous diffusion coefficient. This equilibrium assumption holds when
surface equilibration is fast relative to advection and diffusion over the
resolved transport length scale. Kinetic surface reactions can in principle be
handled by the same interface but are not used in the present work.

The solver uses first-order Lie splitting \citep{carrayrou2004operator}. The
splitting introduces a global error of $\mathcal{O}(\Delta t)$ in the
operator-split solution; the local error over one step is
$\mathcal{O}(\Delta t^2)$. For the timesteps used here, the advective
Courant--Friedrichs--Lewy (CFL) number remains $O(1)$, so the
operator-splitting error is of the same formal order as the first-order
transport discretization. The resulting algorithm keeps the conservative
transport update and the local PHREEQC equilibrium update as distinct
operators within each timestep.

At the algorithmic level, one timestep of size $\Delta t$ is:
\begin{align}
\widetilde{\bm{C}}^{\,n+1} &= \mathcal{T}_{\Delta t}(\bm{C}^{n}, \theta, \bm{u}_D, \bm{D}),\\
(\bm{C}^{n+1}, \bm{\Gamma}^{n+1}) &=
\mathcal{R}_{\Delta t}(\widetilde{\bm{C}}^{\,n+1}, \bm{\Gamma}^{n}, \theta, a_{\mathrm{surf}}),
\end{align}
where $\mathcal{T}_{\Delta t}$ is the conservative finite-volume transport
operator (advancing $\theta\bm{C}$ in DBS form) and $\mathcal{R}_{\Delta t}$
is the local PHREEQC equilibrium operator that updates both aqueous and
surface states simultaneously within each cell.

\subsection{Analytical Dilute Limit}

In the dilute linear-equilibrium limit, surface complexation reduces to a
Henry-type isotherm. The standard Henry coefficient is $K_s = \Gamma / C$
(units: m), the ratio of surface concentration $\Gamma$ (mol\,m$^{-2}$) to
aqueous concentration $C$ (mol\,m$^{-3}$):
\begin{equation}
\Gamma = K_s C.
\end{equation}
In practice, $K_s$ is not prescribed directly but is
extracted numerically from a pair of batch equilibrations at two different
aqueous concentrations:
$K_s = \Delta\Gamma / \Delta C$.
This avoids any ambiguity in the PHREEQC database-to-field unit mapping and
provides the coefficient that enters the retardation formula directly.

Substituting into the inventory Eq.~\eqref{eq:inventory} and taking the
time derivative yields the retarded 1D transport equation
\begin{equation}
\left(
\theta + a_{\mathrm{surf}}K_s
\right)\frac{\partial C}{\partial t}
+ \frac{\partial (u_D C)}{\partial x} = 0.
\end{equation}
The corresponding retardation factor is
\begin{equation}
R = 1 + \frac{a_{\mathrm{surf}}\,K_s}{\theta},
\label{eq:retardation}
\end{equation}
Both $a_{\mathrm{surf}}$ and $\theta$ can vary spatially, so
Eq.~\eqref{eq:retardation} also gives the local retardation at any point in
a heterogeneous domain. The analytical advection-diffusion profile against
which the numerical results are compared follows \citet{ogata1961solution},
with the advection velocity $u_D/\theta$ replaced by $u_D/(\theta R)$
and the diffusivity replaced by $D/R$.

\subsection{Validation Requirements}

The validation programme presented in this paper is designed to establish four
criteria. First, the formulation must reduce to simpler known limits. Second,
the local surface-coupled chemistry must conserve the intended field-level
inventory. Third, the formulation must reproduce the dilute analytical
retardation limit. Fourth, the same formulation must remain applicable without
structural change as the application domain expands from 1D benchmarks to 2D
and 3D image-based cases.
